\documentclass[../../../include/open-logic-section]{subfiles}

\begin{document}

\olfileid{sth}{story}{approach}
\olsection{累积迭代进路}

关于如何将集合分层为阶段，以下是一个稍详细的说明：

\begin{quote}
	集合是在\emph{阶段}中形成的。对于每个阶段 $S$，都存在某些\emph{先于} $S$ 的阶段。在阶段 $S$，每个由在 $S$ 之前的阶段中形成的集合所组成的聚集，都被形成为一个集合。除了在各阶段中形成的集合之外，不存在其他集合。\citep[p.~323]{Shoenfield:AST}
\end{quote}

这就勾勒出了\emph{累积迭代的集合概念}的轮廓。它将构成本书 \olref[sth][][]{part} 所介绍的形式集合论的基础。

让我们来做更详细的探讨。正如 Shoenfield（肖菲尔德）所描述的过程，在每个阶段，我们都从更早阶段可用的集合中形成新的集合。因此，在 Shoenfield 的图景中，在初始阶段（即阶段 $0$），没有任何\emph{更早的}阶段，\emph{因此更不可能}有来自更早阶段的可用集合。\footnote{为何应当假设存在\emph{第一个}阶段？参见 \olref[z][story]{sec} 中关于 \stagesord{} 的脚注。}因此我们只形成一个集合：不含任何元素的集合 $\emptyset$。在阶段 $1$，恰好有一个来自更早阶段的可用集合，因此只形成一个新集合 $\{\emptyset\}$。在阶段 $2$，有两个来自更早阶段的可用集合，我们形成两个新集合 $\{\{\emptyset\}\}$ 和 $\{\emptyset, \{\emptyset\}\}$。在阶段 $3$，有四个来自更早阶段的可用集合，于是我们形成十二个新集合……这样一来，集合的累积迭代图景大致如下（数字表示阶段）：
\begin{center}
	\begin{tikzpicture}[scale=0.6]
	\tikzset{cut_here/.style={densely dotted}}
	\draw (-4,6) -- (-1, 0)-- (2,6);
	\draw[cut_here] (-5,8)--(-4,6);
	\draw[cut_here] (2,6)--(3,8);
	\draw(-1.5, 1)--(-0.5, 1);
	\draw(-2, 2)--(0, 2);
	\draw(-2.5, 3)--(0.5,3);
	\draw(-3, 4)--(1, 4);
	\draw(-3.5, 5)--(1.5, 5);
	\draw(-4, 6)--(2, 6);
	\node[label] at (0, 0) {\small 0};
	\node[label] at (0.5, 1) {\small 1};
	\node[label] at (1, 2) {\small 2};
	\node[label] at (1.5, 3) {\small 3};
	\node[label] at (2, 4) {\small 4};
	\node[label] at (2.5, 5) {\small 5};
	\node[label] at (3, 6) {\small 6};
	\end{tikzpicture}
\end{center}

那么：我们为何要接受这个故事呢？

一个原因是，这是一个令人满意且易于驾驭的故事。鉴于最显而易见的故事（即朴素概括）已经失败，我们正需要某种令人满意的东西。

但这个故事不仅仅是令人满意。我们有充分的理由相信，任何基于这个故事建立的集合论都将是\emph{一致的}。原因如下。

根据累积迭代的集合概念，我们在各个阶段中形成集合；且其元素必须是\emph{已经}可用的对象。因此，对于任何阶段~$S$，我们可以形成集合
\[
	R_S = \Setabs{x}{x \notin x \text{并且 $x$ 在 $S$ 之前可用}}
\]
现在，证明 Russell悖论时所涉及的推理将表明，$R_S$ 本身在阶段 $S$ 之前是不可用的。而这并不矛盾。此外，如果我们接受累积迭代的集合概念，那么我们甚至不应该\emph{期望}能够形成 Russell 集本身。因为那将是所有“终将可用”的非自属集合的集合。简而言之：在累积迭代的故事看来，我们（可证地）不能形成 Russell 集这一事实并不\emph{令人意外}；这正是我们所\emph{预测}的。

%In one sense, then, the cumulative-iterative conception of set yields a response to Russell's Paradox which is rather like the \emph{predicativist}'s response. After all, the predicativist said that the collection of all non-self-membered sets$_n$ is a set$_{n+1}$, and this set$_n$ will not be self-membered. (Indeed, most predicativists will treat it as \emph{ungrammatical} to try to ask whether a set is self-membered.)
%
%But there is an important difference between the cumulative-iterative approach and the predicativist approach: the cumulative-iterative approach treats all of our entities as being \emph{of the same kind}. The predicativist will presumably have an empty set$_0$ (i.e., a set$_0$ with no !!{element}s), and an empty set$_1$ (i.e., a set$_1$ with no !!{element}s), and these will be \emph{different entities}. But on the cumulative-iterative approach, there is just one empty set, $\emptyset$, and it will be available at every stage of the hierarchy.

%Indeed, the Axiom of Extensionality, stated at the start of this chapter, will hold true of the elements in this hierarchy (so that there can be \emph{only one} “empty” set). But I do not intend to use the cumulative-iterative conception to justify Extensionality (see \cite{Boolos1971}). Again: I suggest that we should accept Extensionality, just because we are interested in extensional collections.

\end{document}